\par\indent\SourceLeaf{225}%
III. Der Winkel zwischen zwei Ebenen ist gleich dem sphärischen Winkel zwischen den Großkreisen, die sie darstellen. Er wird daher auch durch den Bogen zwischen den Polen dieser Großkreise gemessen. Ebenso wird die Neigung einer Geraden gegen eine Ebene durch den Bogen gemessen, der vom Punkt, welcher der Richtung der Geraden entspricht, senkrecht zu dem Großkreis gezogen wird, der die Lage der Ebene darstellt.

IV. Bezeichnen \(x,y,z;\ x',y',z'\) die Koordinaten zweier Punkte, \(r\) ihren Abstand und \(L\) den Punkt der Kugeloberfläche, der die Richtung der vom ersten zum zweiten Punkt gezogenen Geraden darstellt, so gilt
\[
x'=x+r\cos(1)L,\qquad y'=y+r\cos(2)L,\qquad z'=z+r\cos(3)L
\]
\par\indent V. Daraus folgt leicht, dass allgemein gilt
\[
\cos(1)L^2+\cos(2)L^2+\cos(3)L^2=1
\]
\noindent und dass für einen beliebigen anderen Punkt \(L'\) der Kugeloberfläche gilt
\[
\cos(1)L\mathord{.}\cos(1)L'+\cos(2)L\mathord{.}\cos(2)L'+\cos(3)L\mathord{.}\cos(3)L'=\cos LL'
\]
\par\indent VI. \textsc{Satz.} \textit{Bezeichnen \(L,L',L'',L'''\) vier Punkte auf der Kugeloberfläche und \(A\) den Winkel, den die Bögen \(LL',L''L'''\) in ihrem Schnittpunkt bilden, so gilt}
\[
\cos LL''\mathord{.}\cos L'L'''-\cos LL'''\mathord{.}\cos L'L''=\sin LL'\mathord{.}\sin L''L'''\mathord{.}\cos A
\]
\par\indent\textit{Beweis.} Der Buchstabe \(A\) bezeichne auch den Schnittpunkt selbst, und wir setzen
\[
AL=t,\quad AL'=t',\quad AL''=t'',\quad AL'''=t'''
\]
\noindent Dann gilt
\[
\begin{aligned}
\cos LL''&=\cos t\cos t''+\sin t\sin t''\cos A\\
\cos L'L'''&=\cos t'\cos t'''+\sin t'\sin t'''\cos A\\
\cos LL'''&=\cos t\cos t'''+\sin t\sin t'''\cos A\\
\cos L'L''&=\cos t'\cos t''+\sin t'\sin t''\cos A
\end{aligned}
\]
\noindent und damit
\[
\begin{aligned}
&\cos LL''\mathord{.}\cos L'L'''-\cos LL'''\mathord{.}\cos L'L''\\
&\quad=\cos A(\cos t\cos t''\sin t'\sin t'''+\cos t'\cos t'''\sin t\sin t''\\
&\qquad\qquad-\cos t\cos t'''\sin t'\sin t''-\cos t'\cos t''\sin t\sin t''')\\
&\quad=\cos A(\cos t\sin t'-\sin t\cos t')(\cos t''\sin t'''-\sin t''\cos t''')\\
&\quad=\cos A\mathord{.}\sin(t'-t)\mathord{.}\sin(t'''-t'')\\
&\quad=\cos A\mathord{.}\sin LL'\mathord{.}\sin L''L'''
\end{aligned}
\]
